% ===== 原文 PDF pp.13--16（印刷 pp.12--15）=====
% tr-01 · tr01b.tex · tag 区间 2.11--2.14（4 条）
% 内容：共形流形与绝对受保护 / 定义 2.1、2.2 / 表 1 / §2.2 标题 / 超 Poincaré 代数 (2.11)--(2.14)

% [原文 p.12]
这一现象表明，短多重态的谱可以在保持超共形对称性的连续形变下发生改变，我们把这类形变所在的流形称为\textbf{共形流形}。当沿着共形流形移动时，一些长多重态可能变成短多重态，而一些短多重态也可能结合成长多重态。因此，短多重态的谱只在复合规则的意义下受保护，这类数据由\textbf{超共形指标} [100] 捕捉。然而，如果某些短多重态从不进入复合规则的右端，它们就真正受到超共形代数的保护，可以沿着共形流形被无歧义地追踪。这些多重态将被称为\textbf{绝对受保护}（absolutely protected）[5]。\footnote{另见关于此类手征算子非重整化性质的早期工作 [101--112]。}值得注意的是，某些 $A$ 型算子也可以是绝对受保护的。对于 $\cN=2$ SCFT 的情形，多重态列于取自 [5] 的表 1 中（另见 [96] 一篇较新的综述）。注意完整的 $\cN=2$ 超共形多重态一般涉及一对手征与反手征缩短条件。
% [Tutor Note: 「绝对受保护」的判据是「从不作为零态出现在复合规则的右端」：$B$ 型多重态因为量子数受严格等式约束、与其他幺正表示隔着有限间隙，天然满足这一点；原文特意提醒 $A$ 型也可能绝对受保护——此时它虽在阈值上，却找不到能与之复合的长多重态（取决于具体理论的算子谱）。因此「短」不等于「受保护」，只有「绝对受保护」才能在共形流形上追踪。]

为后续之用，我们挑出两个重要的短多重态：

\begin{definitionbox}{定义 2.1：库仑分支多重态}
库仑分支多重态是类型为 $\bar LB_1[0,0]^{(0;r)}$（$r>1$）与 $\bar A_2B_1[0,0]^{(0;r)}$（$r=1$，对应自由矢量多重态）的短多重态。其超共形主算子是标度维数为 $\Delta=r$ 的标量算子，称为库仑分支\textbf{手征主算子}，它们生成库仑分支手征环。此外还有反手征主算子，对应手征主算子的共轭表示。
\end{definitionbox}
% [Tutor Note: 两种类型都能在表 1 中定位（表 1 的读法：行 = 手征型 $X$，列 = 反手征型 $\bar Y$，格中即 $X\bar Y$ 多重态）。$\bar LB_1[0,0]^{(0;r)}$ 对应 $B_1$ 行 × $\bar L$ 列：$\bar L$ 列对 $B_1$ 行给出 $r>\frac{\bar j+2}{2}$，代入 $\bar j=0$ 得 $r>1$，且 $\Delta=R+r=r$（$R=0$）；$\bar A_2B_1[0,0]^{(0;r)}$ 对应 $B_1$ 行 × $\bar A_2$ 列：$r=1$、$\Delta=1+R=1$。定义里把反手征型写在前面，与表 1「行在前、列在后」只是书写顺序不同，并非不同的多重态。]

\begin{definitionbox}{定义 2.2：希格斯分支多重态}
希格斯分支多重态是类型为 $B_1\bar B_1[0,0]^{(R;0)}$ 的短多重态。其超共形主算子是权重为 $\Delta=R$、在 $\fu(1)_r$ 下不带荷的标量，但它们在 $\fsu(2)_R$ 的维数为 $R+1$ 的非平凡表示中变换。它们被称为希格斯分支\textbf{手征主算子}，并生成希格斯分支手征环。
\end{definitionbox}
% [Tutor Note: 对照表 1 的 $B_1$ 行 × $\bar B_1$ 列：$[j=0,\bar j=0]^{(R;r=0)}_\Delta$ 且 $\Delta=R$，与定义 2.2 完全一致；这也解释了为何希格斯分支主算子 $\fu(1)_r$ 荷为零而 $\fsu(2)_R$ 表示为 $R+1$ 维（$R\in\Z$ 是 $\fsu(2)_R$ 自旋的两倍，故维数为 $R+1$）。]

这些多重态的一个重要特征是，它们的超共形主算子可以发展出连续的 VEV，从而自发破缺共形对称性、但保留 $\cN=2$ 超 Poincaré 对称性，导致 SCFT 出现一个超对称真空的模空间。视哪一类超共形主算子获得 VEV，模空间相应的分支通常被称为库仑分支、希格斯分支

% [原文 p.13]
% 表 1（原文 p.14 为横排版面；此处用 booktabs 重排，行列归属逐格核对）
\begin{center}
\begingroup
\small
\setlength{\tabcolsep}{4pt}
\resizebox{\textwidth}{!}{%
\begin{tabular}{c cccc}
\toprule
 & $\bar L$ & $\bar A_1$ & $\bar A_2$ & $\bar B_1$\\
\midrule
$L$ &
\begin{tabular}[c]{@{}c@{}}$[j,\bar j]^{(R;r)}_\Delta$\\[2pt]$\Delta>2+R+\max\{j-r,\bar j+r\}$\end{tabular} &
\begin{tabular}[c]{@{}c@{}}$[j,\bar j\ge1]^{(R;\,r<\frac{\bar j-j}{2})}_\Delta$\\[2pt]$\Delta=2+R+\bar j-r$\end{tabular} &
\begin{tabular}[c]{@{}c@{}}$[j,\bar j=0]^{(R;\,r<-\frac{j}{2})}_\Delta$\\[2pt]$\Delta=2+R-r$\end{tabular} &
\begin{tabular}[c]{@{}c@{}}$[j,\bar j=0]^{(R;\,r<-\frac{j+2}{2})}_\Delta$\\[2pt]$\Delta=R-r$\end{tabular}\\
\midrule
$A_1$ &
\begin{tabular}[c]{@{}c@{}}$[j\ge1,\bar j]^{(R;\,r>\frac{\bar j-j}{2})}_\Delta$\\[2pt]$\Delta=2+R+j+r$\end{tabular} &
\begin{tabular}[c]{@{}c@{}}$[j\ge1,\bar j\ge1]^{(R;\,r=\frac{\bar j-j}{2})}_\Delta$\\[2pt]$\Delta=2+R+\frac12(j+\bar j)$\end{tabular} &
\begin{tabular}[c]{@{}c@{}}$[j\ge1,\bar j=0]^{(R;\,r=-\frac{j}{2})}_\Delta$\\[2pt]$\Delta=2+R+\frac12 j$\end{tabular} &
\begin{tabular}[c]{@{}c@{}}$[j\ge1,\bar j=0]^{(R;\,r=-\frac{j+2}{2})}_\Delta$\\[2pt]$\Delta=1+R+\frac12 j$\end{tabular}\\
\midrule
$A_2$ &
\begin{tabular}[c]{@{}c@{}}$[j=0,\bar j]^{(R;\,r>\frac{\bar j}{2})}_\Delta$\\[2pt]$\Delta=2+R+r$\end{tabular} &
\begin{tabular}[c]{@{}c@{}}$[j=0,\bar j\ge1]^{(R;\,r=\frac{\bar j}{2})}_\Delta$\\[2pt]$\Delta=2+R+\frac12\bar j$\end{tabular} &
\begin{tabular}[c]{@{}c@{}}$[j=0,\bar j=0]^{(R;\,r=0)}_\Delta$\\[2pt]$\Delta=2+R$\end{tabular} &
\begin{tabular}[c]{@{}c@{}}$[j=0,\bar j=0]^{(R;\,r=-1)}_\Delta$\\[2pt]$\Delta=1+R$\end{tabular}\\
\midrule
$B_1$ &
\begin{tabular}[c]{@{}c@{}}$[j=0,\bar j]^{(R;\,r>\frac{\bar j+2}{2})}_\Delta$\\[2pt]$\Delta=R+r$\end{tabular} &
\begin{tabular}[c]{@{}c@{}}$[j=0,\bar j\ge1]^{(R;\,r=\frac{\bar j+2}{2})}_\Delta$\\[2pt]$\Delta=1+R+\frac12\bar j$\end{tabular} &
\begin{tabular}[c]{@{}c@{}}$[j=0,\bar j=0]^{(R;\,r=1)}_\Delta$\\[2pt]$\Delta=1+R$\end{tabular} &
\begin{tabular}[c]{@{}c@{}}$[j=0,\bar j=0]^{(R;\,r=0)}_\Delta$\\[2pt]$\Delta=R$\end{tabular}\\
\bottomrule
\end{tabular}}
\endgroup
\end{center}
\tabtitle{1}{幺正的 4d $\cN=2$ 超共形多重态（改编自 [5]，并如 (2.3) 那样对 $\mathrm{U}(1)_r$ 荷作了小幅修改）。}
% [Tutor Note: 这张表要按「行 × 列」读：第 $X$ 行与第 $\bar Y$ 列的交叉格是 $X\bar Y$ 多重态，格中第一行是超共形主算子的量子数与缩短条件，第二行是 $\Delta$ 的幺正性界。行、列各自的缩短类型决定了格中 $(R;r)$ 上标里的等式/不等式，例如 $A_1$ 行要求 $j\ge1$，$A_2$ 行要求 $j=0$。对角线上那个格（$L$ 行 × $\bar L$ 列）不满足任何缩短条件，故只有不等式 $\Delta>2+R+\max\{j-r,\bar j+r\}$，这就是长多重态。]

% [原文 p.14]
与混合分支，这与相应多重态的名字相符。

\SUB{2.2}{4d $\cN=2$ 多重态与一般的 $\cN=2$ 拉格朗日量}

到目前为止，我们关注的是 $\cN=2$ SCFT 中算子的一般性质，这部分内容完全普适、但略显抽象。寻找并研究 CFT 及其算子谱在拉格朗日描述下的实现往往很有启发性。为此，我们在此回顾 $\cN=2$ 拉格朗日理论（不必是共形的）的一般性质。

让我们更仔细地考察 4d $\cN=2$ 超 Poincaré 代数。如前所述，它来自 4d Poincaré 代数 $(P_\mu,M_{\mu\nu})$ 被两个超荷 $\cQ,\tilde\cQ$ 的扩充，相应的量子数为
\begin{equation}
  \left(\cQ^m_\alpha,\tilde\cQ^m_{\dot\alpha}\right)\in
  \left([1,0]^{(1;\frac12)}_{\frac12},\,[0,1]^{(1;-\frac12)}_{\frac12}\right),
  \tag{2.11}
\end{equation}
它们具有如下非平凡的非对易关系（完整的对易与反对易关系见例如 [113--116]）
\begin{equation}
  \begin{aligned}
    \left\{\cQ^m_\alpha,\tilde\cQ^n_{\dot\beta}\right\}&=2(\sigma^\mu)_{\alpha\dot\beta}\,\epsilon^{mn}P_\mu\,,\\
    \left\{\cQ^m_\alpha,\cQ^n_\beta\right\}&=\delta_{\alpha\beta}\,\epsilon^{mn}\mathbf Z\,,
  \end{aligned}
  \tag{2.12}
\end{equation}
其中 $\sigma^\mu=\{1,\sigma^i\}$ 是手征表示下的 4d gamma 矩阵，$\sigma^i$ 是通常的 Pauli 矩阵，$\epsilon^{mn}$ 是 Levi-Civita 张量，$P_\mu$ 是 4d 平移的生成元。更有意思的是 $\mathbf Z$，它被称为\textbf{中心荷}，代表该反对易关系的一个中心扩张。中心扩张在描述 4d $\cN=2$ SCFT 在 IR 的有效场论时将起重要作用。最后要注意的是，在我们的归一化下，由式 (2.3) 中超荷在 $\fu(1)_r$ 下的荷可知，$\mathbf Z$ 将带荷 $+1$。
% [Tutor Note: (2.12) 第二式的荷是自洽的：由 (2.3)，$\cQ$ 的 $\fu(1)_r$ 荷为 $+\frac12$，故 $\{\cQ,\cQ\}$ 的总 $\fu(1)_r$ 荷为 $1$，正是原文所说「$\mathbf Z$ 带荷 $+1$」；又因 $\epsilon^{mn}$ 是 $\SU(2)_R$ 的反对称不变量，$\mathbf Z$ 在 $\SU(2)_R$ 下为单态（这正是「中心」的含义：与所有生成元对易）。另注意 (2.11) 中 $R=1$ 表示 $\fsu(2)_R$ 自旋为 $\frac12$（$R$ 是自旋的两倍），$r=\pm\frac12$ 来自 (2.3)。]

要构造 4d $\cN=2$ 拉格朗日量，我们需要知道上述超对称代数的不可约表示，它们分为两类：\textbf{矢量多重态}与\textbf{超多重态}。
\begin{itemize}
\item \textbf{矢量多重态} $V_{\cN=2}$ 包含一个复标量 $\Phi$ 与两个复 Weyl 规范微子 $\lambda_\alpha,\tilde\lambda_{\dot\alpha}$，以及一个反对称二形式 $F_{\mu\nu}:=\partial_\mu A_\nu+[A_\mu,A_\nu]$，它们都在规范群 $G$ 的伴随表示中。整个多重态可以由标量场 $\Phi$ 构造出来，而 $\Phi$ 进一步被手征超荷湮灭：
% [原文 p.15]
\begin{equation}
  \cQ^m|\Phi\rangle=0\,,\qquad m=1,2\,.
  \tag{2.13}
\end{equation}
然后可以通过用反手征超荷 $\tilde\cQ^n$（$n=1,2$）作用于态 $|\Phi\rangle$，把其他场作为超对称后裔得到：作用一次得到两个规范微子 $\lambda_\alpha,\tilde\lambda_{\dot\alpha}$，作用两次得到场强 $F_{\mu\nu}$。再次说明，上标 $m,n$ 表示 $\SU(2)_R$ 对称性的双重态指标。在我们的记号下，标量场 $\Phi$ 在 $\mathrm{U}(1)_r$ R 对称性下带荷 $1$；并且由超对称代数的反对易关系可见，这两个场 $\lambda^n_\alpha$、$A_\mu$ 分别在 $\mathrm{U}(1)_r$ R 对称性下带荷 $1/2$ 与 $0$。作为自由的 4d $\cN=2$ SCFT，一个 $\cN=2$ 矢量多重态 $V_{\cN=2}$ 实现表 1 中引入的短多重态 $A_2\bar B_1[0,0]^{(0;1)}_1$，其中 $\Phi$ 是超共形主算子。
% [Tutor Note: 原文此处把矢量多重态写成 $A_2\bar B_1[0,0]^{(0;1)}_1$（已用 1800 dpi 裁图复核：横杠确实在 $B_1$ 上，即 $A_2\bar B_1$）。但按表 1，$A_2$ 行 × $\bar B_1$ 列给出 $r=-1$、$\Delta=1+R$，与上标中的 $r=1$、$\Delta=1$ 不符；而能给出 $r=1,\Delta=1+R$ 的是 $B_1$ 行 × $\bar A_2$ 列，即 $\bar A_2B_1$——这正是定义 2.1 中对 $r=1$（自由矢量多重态）所写的类型。故此处疑为横杠标错位置（应为 $\bar A_2B_1$）。译文照录原文，读者按 $\bar A_2B_1[0,0]^{(0;1)}_1$ 理解。]
\item \textbf{超多重态} $H_{\cN=2}$ 涉及两个复标量 $(q,\tilde q)$ 与 Weyl 费米子 $(\psi_\alpha,\tilde\psi_{\dot\alpha})$。这些标量被一半的超荷湮灭，只不过手征性不同，即
\begin{equation}
  \cQ^1|q\rangle=\tilde\cQ^1|q\rangle=\cQ^1|\tilde q\rangle=\tilde\cQ^1|\tilde q\rangle=0\,,
  \tag{2.14}
\end{equation}
其中上标 $1$ 是 $\SU(2)_R$ 双重态指标、其 Cartan 荷为 $+\frac12$，而同样的荷由标量 $(q,\tilde q)$ 携带。现在我们可以用其余超荷作用以得到超对称后裔，它们是费米子。这些标量在 $\mathrm{U}(1)_r$ R 对称性下不带荷，而 $\SU(2)_R$ 把 $(q,(\tilde q)^\dagger)$ 作为双重态作用。费米子在 $\SU(2)_R$ 下不变，并分别在 $\mathrm{U}(1)_r$ 对称性下带 $\pm1/2$ 的荷。超多重态还带有一个额外的 $\SU(2)$ 味对称性，它与 R 对称性对易，且标量与费米子在其中都按双重态变换。在 4d $\cN=2$ 超共形表示的语境下，这样的超多重态实现短多重态 $B_1\bar B_1[0,0]^{(1;0)}_1$。一般地，若有 $n$ 个超多重态，味对称性扩大为 $\mathrm{USp}(2n)$，标量记为 $(q^A,\tilde q_A)$，其中 $A$ 是 $\mathrm{USp}(2n)$ 的 $2n$ 维基本表示的指标。
% [Tutor Note: (2.14) 中 $\cQ^1$ 与 $\tilde\cQ^1$ 同时湮灭 $|q\rangle,|\tilde q\rangle$：这正是超多重态「半缩短」的来源——只保留了 $(\cQ,\tilde\cQ)$ 的一半（上标固定为 $1$）。$\SU(2)_R$ 把 $(q,(\tilde q)^\dagger)$ 组成双重态（而不是 $(q,\tilde q)$），所以 $(\tilde q)^\dagger$ 才带 $-\frac12$ 的 $I_3$ 荷；这也解释了为何原文强调标量带 Cartan 荷 $+\frac12$。]
\end{itemize}

引入 4d $\cN=2$ 超对称多重态之后，我们现在可以写出 4d $\cN=2$ 拉格朗日量的一般形式。超对称对拉格朗日量的形式施加强烈的限制，并且存在一种系统而高效的办法，用 4d $\cN=1$ 的超空间形式来构造由上述超多重态构成的 4d $\cN=2$ 拉格朗日量。在该形式中，通过把通常的时空扩充为带有额外 Grassmann 坐标 $(\theta_\alpha,\bar\theta_{\dot\alpha})$ 的超空间，我们可以构造各种 4d $\cN=1$ 超场。
